% Options for packages loaded elsewhere
\PassOptionsToPackage{unicode}{hyperref}
\PassOptionsToPackage{hyphens}{url}
\PassOptionsToPackage{dvipsnames,svgnames,x11names}{xcolor}
%
\documentclass[
  11pt,
]{article}
\usepackage{amsmath,amssymb}
\usepackage{iftex}
\ifPDFTeX
  \usepackage[T1]{fontenc}
  \usepackage[utf8]{inputenc}
  \usepackage{textcomp} % provide euro and other symbols
\else % if luatex or xetex
  \usepackage{unicode-math} % this also loads fontspec
  \defaultfontfeatures{Scale=MatchLowercase}
  \defaultfontfeatures[\rmfamily]{Ligatures=TeX,Scale=1}
\fi
\usepackage{lmodern}
\ifPDFTeX\else
  % xetex/luatex font selection
\fi
% Use upquote if available, for straight quotes in verbatim environments
\IfFileExists{upquote.sty}{\usepackage{upquote}}{}
\IfFileExists{microtype.sty}{% use microtype if available
  \usepackage[]{microtype}
  \UseMicrotypeSet[protrusion]{basicmath} % disable protrusion for tt fonts
}{}
\makeatletter
\@ifundefined{KOMAClassName}{% if non-KOMA class
  \IfFileExists{parskip.sty}{%
    \usepackage{parskip}
  }{% else
    \setlength{\parindent}{0pt}
    \setlength{\parskip}{6pt plus 2pt minus 1pt}}
}{% if KOMA class
  \KOMAoptions{parskip=half}}
\makeatother
\usepackage{xcolor}
\usepackage[margin=2.2cm]{geometry}
\usepackage{longtable,booktabs,array}
\usepackage{pdflscape}
\usepackage{ragged2e}
\usepackage{calc} % for calculating minipage widths
% Correct order of tables after \paragraph or \subparagraph
\usepackage{etoolbox}
\makeatletter
\patchcmd\longtable{\par}{\if@noskipsec\mbox{}\fi\par}{}{}
\makeatother
% Allow footnotes in longtable head/foot
\IfFileExists{footnotehyper.sty}{\usepackage{footnotehyper}}{\usepackage{footnote}}
\makesavenoteenv{longtable}
\setlength{\emergencystretch}{3em} % prevent overfull lines
\providecommand{\tightlist}{%
  \setlength{\itemsep}{0pt}\setlength{\parskip}{0pt}}
\setcounter{secnumdepth}{-\maxdimen} % remove section numbering
\ifLuaTeX
  \usepackage{selnolig}  % disable illegal ligatures
\fi
\IfFileExists{bookmark.sty}{\usepackage{bookmark}}{\usepackage{hyperref}}
\IfFileExists{xurl.sty}{\usepackage{xurl}}{} % add URL line breaks if available
\urlstyle{same}
\hypersetup{
  colorlinks=true,
  linkcolor={Maroon},
  filecolor={Maroon},
  citecolor={Blue},
  urlcolor={Blue},
  pdfcreator={LaTeX via pandoc}}

\author{}
\date{}

\begin{document}

\hypertarget{la-meta-para-el-examen-capuxedtulos-27-al-final}{%
\section{La Meta para el examen: capítulos 27 al
final}\label{la-meta-para-el-examen-capuxedtulos-27-al-final}}

Documento preparado para estudiar la parte de \emph{La Meta} que entra
en el examen de Gestión de Operaciones. Está escrito para poder cargarlo
en NotebookLM y generar un resumen en audio.

\hypertarget{alcance-seguxfan-los-temarios}{%
\subsection{Alcance según los
temarios}\label{alcance-seguxfan-los-temarios}}

Según \texttt{temario\_i1.md}, en la I1 entraron los capítulos 1 al 15.
Según \texttt{temario\_i2.md}, en la I2 entraron los capítulos 16 al 26.
Según \texttt{temario\_examen.md}, para el examen entra \emph{La Meta}
desde el capítulo 27 inclusive hasta el final.

Por lo tanto, este documento se enfoca en:

\begin{itemize}
\tightlist
\item
  Capítulos que entran al examen: 27 al 40.
\item
  Capítulos que se excluyen como foco principal: 1 al 26.
\item
  Conceptos anteriores que igual se usan como base: throughput,
  inventario, gasto operativo, cuello de botella, eventos dependientes,
  fluctuaciones estadísticas, DBR y los cinco pasos de mejora continua.
\end{itemize}

No es una transcripción del libro. Es una guía de estudio con resumen,
interpretación operacional y posibles afirmaciones Verdadero/Falso.

\hypertarget{resumen-ejecutivo-para-audio}{%
\subsection{Resumen ejecutivo para
audio}\label{resumen-ejecutivo-para-audio}}

En los capítulos 27 al 40, la historia deja de ser solamente la
salvación de una planta y pasa a ser una reflexión sobre cómo
administrar un sistema completo. Alex ya logró mejorar la planta usando
la Teoría de Restricciones, pero la presión no termina: Bill Peach exige
una mejora adicional, aparece la necesidad de vender más, se reducen
tamaños de lote, se cuestionan los indicadores contables tradicionales,
la planta gana flexibilidad, Alex es promovido y el equipo intenta
generalizar lo aprendido a toda la división.

La idea central de esta parte es que la mejora operacional no consiste
en apagar incendios ni en mantener todos los recursos ocupados. Consiste
en entender el sistema, identificar su restricción, explotarla,
subordinar las demás decisiones, elevar la restricción si hace falta y
volver a empezar cuando la restricción cambia. Al final, Alex descubre
que administrar no es memorizar soluciones, sino saber formular las
preguntas correctas: qué cambiar, hacia qué cambiar y cómo provocar el
cambio.

Para el examen, esta parte puede aparecer en Verdadero/Falso,
especialmente en afirmaciones sobre lotes pequeños, lead time,
inventario, throughput, contabilidad tradicional, inercia
organizacional, restricción de mercado, contratos grandes, capacidad
ociosa y procesos de pensamiento.

\hypertarget{personajes-importantes-en-esta-parte}{%
\subsection{Personajes importantes en esta
parte}\label{personajes-importantes-en-esta-parte}}

\hypertarget{alex-rogo}{%
\subsubsection{Alex Rogo}\label{alex-rogo}}

Gerente de planta que ya aplicó principios de la Teoría de Restricciones
y ahora debe demostrar que la mejora no fue casual. En estos capítulos
pasa de administrar una planta a tener responsabilidad divisional. Su
aprendizaje principal es pasar de depender de Jonah a desarrollar su
propio método de pensamiento.

\hypertarget{jonah}{%
\subsubsection{Jonah}\label{jonah}}

Mentor de Alex. No entrega respuestas cerradas, sino preguntas. Su rol
es empujar a Alex a deducir principios generales. En esta parte aparece
menos, porque Alex debe volverse autosuficiente.

\hypertarget{bill-peach}{%
\subsubsection{Bill Peach}\label{bill-peach}}

Ejecutivo superior. Presiona a Alex con metas agresivas y representa la
mirada corporativa que exige resultados financieros. Aunque reconoce
mejoras, sigue dudando de que sean sostenibles.

\hypertarget{lou}{%
\subsubsection{Lou}\label{lou}}

Contador de la planta. Es clave porque muestra que los indicadores
contables tradicionales pueden ocultar la mejora real. Ayuda a traducir
throughput, inventario y gasto operativo a lenguaje financiero.

\hypertarget{bob-donovan}{%
\subsubsection{Bob Donovan}\label{bob-donovan}}

Producción. Representa la preocupación práctica por ejecutar los cambios
en planta, mover recursos y cumplir pedidos.

\hypertarget{stacey}{%
\subsubsection{Stacey}\label{stacey}}

Inventarios y materiales. Es clave para entender cómo la liberación de
material, los buffers y las prioridades afectan WIP y cumplimiento.

\hypertarget{ralph}{%
\subsubsection{Ralph}\label{ralph}}

Datos y sistemas. Ayuda a transformar las ideas de flujo en reglas
operacionales y programación de liberación de materiales.

\hypertarget{johnny-jons}{%
\subsubsection{Johnny Jons}\label{johnny-jons}}

Ventas. Se vuelve importante cuando la restricción deja de ser interna y
aparece la necesidad de conseguir más demanda para usar la capacidad
liberada.

\hypertarget{julie}{%
\subsubsection{Julie}\label{julie}}

Esposa de Alex. En esta parte la historia personal sirve como espejo del
aprendizaje administrativo: también en la vida personal hay que
cuestionar supuestos, metas y reglas heredadas.

\hypertarget{conceptos-base-que-debes-recordar-antes-del-capuxedtulo-27}{%
\subsection{Conceptos base que debes recordar antes del capítulo
27}\label{conceptos-base-que-debes-recordar-antes-del-capuxedtulo-27}}

\hypertarget{la-meta-de-una-empresa}{%
\subsubsection{La meta de una empresa}\label{la-meta-de-una-empresa}}

La meta de una empresa con fines de lucro es ganar dinero ahora y en el
futuro. No basta con producir mucho, ocupar máquinas o reducir costos
localmente. Una acción es productiva si acerca a la empresa a ganar más
dinero de forma sostenible.

\hypertarget{tres-medidas-operacionales}{%
\subsubsection{Tres medidas
operacionales}\label{tres-medidas-operacionales}}

Throughput, inventario y gasto operativo son las tres medidas que
ordenan la toma de decisiones.

Throughput es la tasa a la que el sistema genera dinero mediante ventas.
No es simplemente producir, porque producir algo que no se vende solo
aumenta inventario.

Inventario es el dinero atrapado en el sistema en cosas que se pretende
vender. Incluye materia prima, trabajo en proceso y producto terminado.

Gasto operativo es el dinero que se gasta para transformar inventario en
throughput.

\hypertarget{regla-de-cuello-de-botella}{%
\subsubsection{Regla de cuello de
botella}\label{regla-de-cuello-de-botella}}

Un cuello de botella es un recurso cuya capacidad limita el desempeño
del sistema. Si se pierde una hora en el cuello de botella, se pierde
una hora de capacidad del sistema completo. Si se ahorra una hora en un
recurso que no es cuello de botella, esa hora no necesariamente aumenta
throughput.

\hypertarget{tambor-amortiguador-y-cuerda}{%
\subsubsection{Tambor, amortiguador y
cuerda}\label{tambor-amortiguador-y-cuerda}}

El tambor es el cuello de botella, porque marca el ritmo del sistema. El
amortiguador es protección antes del cuello de botella para que no quede
sin trabajo. La cuerda es la regla que controla la liberación de
material de acuerdo con el ritmo del cuello de botella.

\hypertarget{cinco-pasos-de-mejora-continua}{%
\subsubsection{Cinco pasos de mejora
continua}\label{cinco-pasos-de-mejora-continua}}

\begin{enumerate}
\def\labelenumi{\arabic{enumi}.}
\tightlist
\item
  Identificar la restricción del sistema.
\item
  Explotar la restricción.
\item
  Subordinar todo lo demás a la decisión anterior.
\item
  Elevar la restricción.
\item
  Si la restricción se rompe, volver al paso 1 y evitar que la inercia
  se transforme en la nueva restricción.
\end{enumerate}

\hypertarget{capuxedtulo-27-la-presiuxf3n-por-mejorar-otro-15}{%
\subsection{Capítulo 27: La presión por mejorar otro
15\%}\label{capuxedtulo-27-la-presiuxf3n-por-mejorar-otro-15}}

\hypertarget{quuxe9-ocurre}{%
\subsubsection{Qué ocurre}\label{quuxe9-ocurre}}

La planta de Alex muestra mejoras reales: mejor cumplimiento, menos
atrasos y mayor control del flujo. Sin embargo, en la reunión
corporativa la reacción no es una celebración total. La división
completa sigue en problemas, y Bill Peach no quiere decidir solo por un
buen mes. Le exige a Alex una mejora adicional importante,
aproximadamente un 15\%, para justificar mantener la planta y confiar en
que el cambio es sostenible.

Alex acepta el desafío, aunque internamente entiende que la mejora
adicional no se resuelve solo produciendo mejor. Si la planta ya aumentó
su capacidad efectiva y redujo atrasos, entonces el nuevo problema puede
estar en vender más. Esto es clave: cuando se mejora una restricción
interna, la restricción puede moverse al mercado.

En paralelo, Alex intenta reconstruir su relación con Julie. La
conversación personal refuerza una idea del libro: las reglas asumidas
deben revisarse. Así como la planta tenía supuestos falsos sobre
eficiencia, el matrimonio también tiene supuestos no discutidos sobre
qué significa funcionar bien.

\hypertarget{concepto-operacional}{%
\subsubsection{Concepto operacional}\label{concepto-operacional}}

La restricción del sistema puede cambiar. Al comienzo, la restricción
estaba dentro de la planta: máquinas, tratamiento térmico, programación
y flujo. Después de mejorar esos puntos, aparece la necesidad de
demanda. La empresa puede tener capacidad disponible, pero si no hay
ventas, el throughput no aumenta.

\hypertarget{lo-evaluable}{%
\subsubsection{Lo evaluable}\label{lo-evaluable}}

El capítulo muestra que mejorar operaciones no basta si el mercado no
absorbe la capacidad. También muestra que una mejora mensual aislada no
garantiza cambio estructural.

\hypertarget{posible-vf}{%
\subsubsection{Posible V/F}\label{posible-vf}}

\textbf{Afirmación:} Después de mejorar la planta, Alex ya no necesita
preocuparse por ventas porque el problema era puramente productivo.\\
\textbf{Respuesta:} Falso. Al liberar capacidad interna, la restricción
puede moverse al mercado. Para aumentar throughput se necesitan ventas
reales, no solo capacidad productiva.

\hypertarget{capuxedtulo-28-reducir-lotes-para-bajar-inventario-y-lead-time}{%
\subsection{Capítulo 28: Reducir lotes para bajar inventario y lead
time}\label{capuxedtulo-28-reducir-lotes-para-bajar-inventario-y-lead-time}}

\hypertarget{quuxe9-ocurre-1}{%
\subsubsection{Qué ocurre}\label{quuxe9-ocurre-1}}

Alex consulta a Jonah sobre cómo lograr la mejora adicional. Jonah
propone reducir los tamaños de lote, especialmente en recursos no cuello
de botella. La idea parece contraintuitiva desde la mirada tradicional,
porque más lotes pueden implicar más setups. Sin embargo, si el recurso
no es cuello de botella, parte de su capacidad ociosa puede usarse para
más preparaciones sin afectar el throughput total.

Jonah también ordena el tiempo que una pieza pasa en la planta en cuatro
componentes:

\begin{itemize}
\tightlist
\item
  Tiempo de preparación o setup.
\item
  Tiempo de proceso.
\item
  Tiempo de cola.
\item
  Tiempo de espera por otras piezas o por ensamble.
\end{itemize}

La reducción de lotes disminuye el trabajo en proceso y acorta el tiempo
total de flujo. Esto permite prometer plazos de entrega menores. Alex lo
conecta con ventas: si la planta puede entregar rápido y confiable,
puede ganar pedidos que antes no podía aceptar.

\hypertarget{concepto-operacional-1}{%
\subsubsection{Concepto operacional}\label{concepto-operacional-1}}

Reducir lotes de transferencia y proceso puede reducir WIP y lead time.
El impacto no se evalúa solo por eficiencia local o setups, sino por
efecto en throughput, inventario y gasto operativo.

En términos de Ley de Little:

WIP = Throughput x Tiempo de Flujo.

Si el throughput se mantiene y baja el WIP, baja el tiempo de flujo. Si
baja el tiempo de flujo, la planta gana velocidad de respuesta
comercial.

\hypertarget{lo-evaluable-1}{%
\subsubsection{Lo evaluable}\label{lo-evaluable-1}}

El capítulo es muy preguntable. La trampa típica es decir que lotes más
pequeños siempre son peores porque aumentan setups. La respuesta
correcta depende de si el recurso es cuello de botella o no.

\hypertarget{posible-vf-1}{%
\subsubsection{Posible V/F}\label{posible-vf-1}}

\textbf{Afirmación:} Reducir el tamaño de lote siempre reduce la
eficiencia y por eso empeora el desempeño global.\\
\textbf{Respuesta:} Falso. Puede bajar la eficiencia local, pero también
reduce WIP y lead time. Si se hace en recursos no cuello de botella,
puede mejorar la respuesta al mercado sin reducir throughput.

\hypertarget{capuxedtulo-29-la-oportunidad-comercial-y-el-pedido-grande}{%
\subsection{Capítulo 29: La oportunidad comercial y el pedido
grande}\label{capuxedtulo-29-la-oportunidad-comercial-y-el-pedido-grande}}

\hypertarget{quuxe9-ocurre-2}{%
\subsubsection{Qué ocurre}\label{quuxe9-ocurre-2}}

Aparece una oportunidad comercial importante: un cliente necesita una
orden grande en poco tiempo. Al principio parece imposible cumplirla sin
desordenar la planta. Sin embargo, el equipo usa lo aprendido: no acepta
el problema tal como viene, sino que cuestiona supuestos.

La solución no es producir todo de una vez. El equipo negocia entregas
parciales y escalonadas. Esto reduce presión sobre la planta, mejora la
factibilidad y satisface mejor al cliente. Los lotes más pequeños y la
programación más flexible permiten responder a una demanda que antes se
habría rechazado.

También aparece una tensión contable: aunque operacionalmente la planta
está mejor, algunos reportes tradicionales pueden verse peor. Esto
ocurre porque la contabilidad de costos puede castigar más setups o
menor eficiencia local, aunque el sistema esté generando más throughput.

\hypertarget{concepto-operacional-2}{%
\subsubsection{Concepto operacional}\label{concepto-operacional-2}}

La flexibilidad tiene valor comercial. Una planta con menor lead time
puede vender más, aunque sus métricas locales no parezcan perfectas.

Este capítulo conecta operaciones con ventas. La mejora de flujo no es
solo interna; permite prometer mejor al cliente, aceptar pedidos
atractivos y convertir capacidad disponible en throughput.

\hypertarget{lo-evaluable-2}{%
\subsubsection{Lo evaluable}\label{lo-evaluable-2}}

Es clave distinguir entre:

\begin{itemize}
\tightlist
\item
  Producir por eficiencia local.
\item
  Producir según necesidad real del mercado.
\end{itemize}

La segunda opción es coherente con TOC. La primera puede crear
inventario inútil.

\hypertarget{posible-vf-2}{%
\subsubsection{Posible V/F}\label{posible-vf-2}}

\textbf{Afirmación:} Aceptar un pedido grande siempre obliga a
producirlo en un único lote grande.\\
\textbf{Respuesta:} Falso. Se pueden negociar entregas parciales. Eso
reduce WIP, mejora flujo y permite cumplir sin saturar innecesariamente
el sistema.

\hypertarget{capuxedtulo-30-el-cliente-satisfecho-y-la-tensiuxf3n-con-el-15}{%
\subsection{Capítulo 30: El cliente satisfecho y la tensión con el
15\%}\label{capuxedtulo-30-el-cliente-satisfecho-y-la-tensiuxf3n-con-el-15}}

\hypertarget{quuxe9-ocurre-3}{%
\subsubsection{Qué ocurre}\label{quuxe9-ocurre-3}}

El cliente importante queda satisfecho con el desempeño de la planta.
Esto valida la idea de que la velocidad de respuesta puede transformarse
en ventaja competitiva. La planta no solo produce mejor; ahora puede
vender una promesa distinta: entregar rápido y de forma confiable.

El resultado operacional mejora, pero Alex sigue preocupado porque el
objetivo corporativo exacto puede no haberse alcanzado según los
reportes tradicionales. La planta parece estar cerca, pero no
necesariamente llega al 15\% pedido por Peach bajo las métricas usuales.

\hypertarget{concepto-operacional-3}{%
\subsubsection{Concepto operacional}\label{concepto-operacional-3}}

Una mejora real puede no verse correctamente si el sistema de medición
está mal diseñado. Si la contabilidad tradicional castiga decisiones que
aumentan throughput, entonces el indicador puede inducir malas
decisiones.

\hypertarget{lo-evaluable-3}{%
\subsubsection{Lo evaluable}\label{lo-evaluable-3}}

La contabilidad tradicional puede entrar en conflicto con TOC. TOC
prioriza throughput global, inventario y gasto operativo, mientras que
la contabilidad de costos puede enfatizar eficiencia local, absorción de
costos y utilización.

\hypertarget{posible-vf-3}{%
\subsubsection{Posible V/F}\label{posible-vf-3}}

\textbf{Afirmación:} Si una mejora operacional no se ve inmediatamente
en los indicadores contables tradicionales, entonces la mejora no
existe.\\
\textbf{Respuesta:} Falso. Los indicadores pueden estar mal alineados.
La mejora debe evaluarse por efecto en throughput, inventario, gasto
operativo y cumplimiento real.

\hypertarget{capuxedtulo-31-la-contabilidad-muestra-otra-realidad-y-alex-es-promovido}{%
\subsection{Capítulo 31: La contabilidad muestra otra realidad y Alex es
promovido}\label{capuxedtulo-31-la-contabilidad-muestra-otra-realidad-y-alex-es-promovido}}

\hypertarget{quuxe9-ocurre-4}{%
\subsubsection{Qué ocurre}\label{quuxe9-ocurre-4}}

Lou ayuda a mostrar que el desempeño real de la planta es mejor que lo
que indicaban ciertos reportes. Al reinterpretar los resultados desde
una lógica más cercana al throughput, se entiende que la mejora de Alex
supera lo que parecía inicialmente.

La planta no se cierra. Además, Bill Peach es promovido y Alex también:
pasa a ocupar un cargo superior en la división. Esto cambia la escala
del problema. Alex ya no debe salvar una planta, sino aplicar el
aprendizaje a varias unidades.

\hypertarget{concepto-operacional-4}{%
\subsubsection{Concepto operacional}\label{concepto-operacional-4}}

El cuello de botella puede dejar de ser técnico y pasar a ser
administrativo. Cuando Alex es promovido, su problema ya no es solo una
máquina o un flujo de materiales; ahora debe cambiar políticas,
métricas, incentivos y formas de pensar.

\hypertarget{lo-evaluable-4}{%
\subsubsection{Lo evaluable}\label{lo-evaluable-4}}

La promoción de Alex muestra que TOC no es una receta local de planta,
sino un método de gestión sistémica.

\hypertarget{posible-vf-4}{%
\subsubsection{Posible V/F}\label{posible-vf-4}}

\textbf{Afirmación:} La Teoría de Restricciones solo sirve para máquinas
y procesos físicos.\\
\textbf{Respuesta:} Falso. En los capítulos finales se generaliza a la
administración de una división, incluyendo políticas, métricas, ventas y
cambio organizacional.

\hypertarget{capuxedtulo-32-de-depender-de-jonah-a-aprender-a-pensar}{%
\subsection{Capítulo 32: De depender de Jonah a aprender a
pensar}\label{capuxedtulo-32-de-depender-de-jonah-a-aprender-a-pensar}}

\hypertarget{quuxe9-ocurre-5}{%
\subsubsection{Qué ocurre}\label{quuxe9-ocurre-5}}

Alex y Julie celebran la promoción. Alex reflexiona sobre Jonah y se
pregunta por qué necesitó tanta ayuda externa para descubrir cosas que,
una vez entendidas, parecen lógicas. Esto lo lleva a una inquietud
central: ¿cómo puede replicar el aprendizaje sin que Jonah esté siempre
presente?

El problema ahora es pedagógico y organizacional. Alex debe lograr que
otros equipos acepten principios que su planta aprendió mediante
experiencia, presión y errores. No basta con ordenar que usen TOC; deben
entender por qué.

\hypertarget{concepto-operacional-5}{%
\subsubsection{Concepto operacional}\label{concepto-operacional-5}}

La implementación de mejoras requiere aprendizaje organizacional. Un
método impuesto sin comprensión puede generar resistencia o aplicación
mecánica.

\hypertarget{lo-evaluable-5}{%
\subsubsection{Lo evaluable}\label{lo-evaluable-5}}

El método socrático de Jonah importa porque reduce resistencia al
cambio. Las personas aceptan mejor una solución cuando participan en
descubrirla.

\hypertarget{posible-vf-5}{%
\subsubsection{Posible V/F}\label{posible-vf-5}}

\textbf{Afirmación:} Jonah es útil porque entrega respuestas operativas
exactas que Alex solo debe copiar.\\
\textbf{Respuesta:} Falso. Jonah guía mediante preguntas para que Alex y
su equipo descubran relaciones causales y aprendan a pensar
sistémicamente.

\hypertarget{capuxedtulo-33-reorganizaciuxf3n-del-equipo}{%
\subsection{Capítulo 33: Reorganización del
equipo}\label{capuxedtulo-33-reorganizaciuxf3n-del-equipo}}

\hypertarget{quuxe9-ocurre-6}{%
\subsubsection{Qué ocurre}\label{quuxe9-ocurre-6}}

Alex asume su nuevo rol y promueve a miembros clave de su equipo. Bob,
Lou, Stacey y Ralph pasan a posiciones donde pueden ayudar a extender el
aprendizaje a la división.

Esto es importante porque el cambio no puede depender de una sola
persona. Alex necesita un equipo capaz de mirar producción, finanzas,
inventario y datos bajo la misma lógica sistémica.

\hypertarget{concepto-operacional-6}{%
\subsubsection{Concepto operacional}\label{concepto-operacional-6}}

La mejora continua necesita estructura organizacional. Si las personas
que entienden el sistema no tienen autoridad, el cambio queda atrapado
en la planta original.

\hypertarget{lo-evaluable-6}{%
\subsubsection{Lo evaluable}\label{lo-evaluable-6}}

TOC requiere coherencia entre áreas. Producción, contabilidad, ventas,
inventario y sistemas deben mirar la misma meta.

\hypertarget{posible-vf-6}{%
\subsubsection{Posible V/F}\label{posible-vf-6}}

\textbf{Afirmación:} Una vez que se identifica la restricción, basta con
que producción actúe; finanzas, ventas y sistemas son secundarios.\\
\textbf{Respuesta:} Falso. La restricción y sus efectos cruzan áreas.
Lou, Johnny, Stacey, Bob y Ralph cumplen roles complementarios en la
mejora.

\hypertarget{capuxedtulo-34-el-desafuxedo-de-escalar-el-muxe9todo}{%
\subsection{Capítulo 34: El desafío de escalar el
método}\label{capuxedtulo-34-el-desafuxedo-de-escalar-el-muxe9todo}}

\hypertarget{quuxe9-ocurre-7}{%
\subsubsection{Qué ocurre}\label{quuxe9-ocurre-7}}

El equipo comienza a trabajar en cómo aplicar el método a toda la
división. Ya no basta con recordar acciones específicas como etiquetas
rojas, buffers o reducción de lotes. Esas fueron soluciones para una
situación concreta. Lo que necesitan ahora es extraer el principio
general.

El desafío es transformar una experiencia exitosa en un método
replicable.

\hypertarget{concepto-operacional-7}{%
\subsubsection{Concepto operacional}\label{concepto-operacional-7}}

No hay que confundir herramienta con principio. Las etiquetas, los
buffers o el programa de Ralph son herramientas. El principio es
subordinar el sistema a su restricción para aumentar throughput con
menos inventario y gasto operativo controlado.

\hypertarget{lo-evaluable-7}{%
\subsubsection{Lo evaluable}\label{lo-evaluable-7}}

Puede aparecer una trampa de copiar soluciones sin diagnosticar. Por
ejemplo: decir que toda planta debe usar exactamente el mismo buffer o
el mismo sistema de etiquetas. Eso es falso: primero se identifica la
restricción del sistema particular.

\hypertarget{posible-vf-7}{%
\subsubsection{Posible V/F}\label{posible-vf-7}}

\textbf{Afirmación:} Como las etiquetas funcionaron en la planta de
Alex, toda división debe usar el mismo sistema sin análisis adicional.\\
\textbf{Respuesta:} Falso. Las herramientas dependen de la restricción y
del sistema. Lo que se replica es el método de análisis, no
necesariamente la herramienta exacta.

\hypertarget{capuxedtulo-35-buscar-el-orden-detruxe1s-del-caos}{%
\subsection{Capítulo 35: Buscar el orden detrás del
caos}\label{capuxedtulo-35-buscar-el-orden-detruxe1s-del-caos}}

\hypertarget{quuxe9-ocurre-8}{%
\subsubsection{Qué ocurre}\label{quuxe9-ocurre-8}}

El equipo reflexiona sobre cómo los científicos clasifican fenómenos
complejos. La analogía es que la división parece caótica, pero puede
existir una estructura lógica subyacente. Si se descubre esa estructura,
los problemas dejan de ser una lista infinita de síntomas y se
transforman en relaciones causa-efecto.

Esta parte prepara la transición desde la gestión de operaciones hacia
los procesos de pensamiento.

\hypertarget{concepto-operacional-8}{%
\subsubsection{Concepto operacional}\label{concepto-operacional-8}}

Un sistema complejo no se administra persiguiendo cada síntoma. Se
administra identificando relaciones causales, restricciones y políticas
que generan muchos efectos a la vez.

\hypertarget{lo-evaluable-8}{%
\subsubsection{Lo evaluable}\label{lo-evaluable-8}}

La mejora continua no es solo técnica; requiere pensamiento lógico. El
gerente debe saber separar síntomas de causas raíz.

\hypertarget{posible-vf-8}{%
\subsubsection{Posible V/F}\label{posible-vf-8}}

\textbf{Afirmación:} Si hay muchos problemas en una organización, cada
problema debe resolverse de forma independiente.\\
\textbf{Respuesta:} Falso. Muchos problemas pueden tener una causa común
o una restricción compartida. TOC busca la causa que gobierna el
desempeño global.

\hypertarget{capuxedtulo-36-formalizaciuxf3n-de-los-cinco-pasos}{%
\subsection{Capítulo 36: Formalización de los cinco
pasos}\label{capuxedtulo-36-formalizaciuxf3n-de-los-cinco-pasos}}

\hypertarget{quuxe9-ocurre-9}{%
\subsubsection{Qué ocurre}\label{quuxe9-ocurre-9}}

El equipo reconstruye lo que hizo en la planta y lo transforma en los
cinco pasos de mejora continua:

\begin{enumerate}
\def\labelenumi{\arabic{enumi}.}
\tightlist
\item
  Identificar la restricción.
\item
  Explotarla.
\item
  Subordinar todo lo demás.
\item
  Elevar la restricción.
\item
  Volver al paso 1 si la restricción se rompió.
\end{enumerate}

Lo importante es que estos pasos no aparecen como teoría abstracta;
nacen de la experiencia concreta de la planta.

\hypertarget{concepto-operacional-9}{%
\subsubsection{Concepto operacional}\label{concepto-operacional-9}}

El Proceso de Mejora Continua, o POOGI, es un algoritmo de gestión. Su
valor está en el orden. No se debe invertir capacidad antes de explotar
la restricción; no se debe optimizar un no cuello de botella antes de
subordinarlo; no se debe detener el proceso cuando una restricción
cambia.

\hypertarget{lo-evaluable-9}{%
\subsubsection{Lo evaluable}\label{lo-evaluable-9}}

El orden de los cinco pasos puede preguntarse directamente. También
pueden aparecer afirmaciones falsas que mezclan explotar con elevar.

\hypertarget{posible-vf-9}{%
\subsubsection{Posible V/F}\label{posible-vf-9}}

\textbf{Afirmación:} Explotar una restricción significa comprar más
capacidad para ella.\\
\textbf{Respuesta:} Falso. Explotar es sacar el máximo de la restricción
con los recursos actuales. Comprar o agregar capacidad corresponde a
elevar.

\hypertarget{capuxedtulo-37-inercia-uxf3rdenes-ficticias-y-capacidad-liberada}{%
\subsection{Capítulo 37: Inercia, órdenes ficticias y capacidad
liberada}\label{capuxedtulo-37-inercia-uxf3rdenes-ficticias-y-capacidad-liberada}}

\hypertarget{quuxe9-ocurre-10}{%
\subsubsection{Qué ocurre}\label{quuxe9-ocurre-10}}

El equipo revisa los cinco pasos y detecta un problema importante: la
inercia. Una política que fue correcta para una restricción anterior
puede transformarse en una restricción nueva si se mantiene sin revisar.

También descubren que se habían generado órdenes ficticias para mantener
ocupados los cuellos de botella. Al eliminar esas órdenes, se libera
capacidad que puede usarse para pedidos reales.

Stacey identifica que el sistema de prioridades debe ajustarse. Johnny
debe buscar mercado para usar capacidad disponible.

\hypertarget{concepto-operacional-10}{%
\subsubsection{Concepto operacional}\label{concepto-operacional-10}}

La inercia es peligrosa porque convierte una solución antigua en una
regla obsoleta. En TOC, cuando una restricción se rompe, hay que volver
al paso 1. Si no se hace, la organización sigue actuando como si el
cuello de botella antiguo todavía gobernara el sistema.

Las órdenes ficticias son un ejemplo de mala subordinación: mantener
ocupado un recurso no sirve si ese trabajo no genera ventas reales.

\hypertarget{lo-evaluable-10}{%
\subsubsection{Lo evaluable}\label{lo-evaluable-10}}

Este capítulo es muy preguntable por la frase del paso 5: no permitir
que la inercia se transforme en restricción.

\hypertarget{posible-vf-10}{%
\subsubsection{Posible V/F}\label{posible-vf-10}}

\textbf{Afirmación:} Si una política funcionó para resolver una
restricción, debe mantenerse permanentemente.\\
\textbf{Respuesta:} Falso. Cuando la restricción cambia, una política
antigua puede convertirse en la nueva restricción. Hay que volver al
paso 1.

\hypertarget{capuxedtulo-38-la-restricciuxf3n-de-mercado-y-el-pedido-europeo}{%
\subsection{Capítulo 38: La restricción de mercado y el pedido
europeo}\label{capuxedtulo-38-la-restricciuxf3n-de-mercado-y-el-pedido-europeo}}

\hypertarget{quuxe9-ocurre-11}{%
\subsubsection{Qué ocurre}\label{quuxe9-ocurre-11}}

Johnny encuentra una oportunidad con un cliente europeo. El precio
exigido es menor que el precio del mercado local, lo que al principio
parece poco atractivo. Sin embargo, Alex analiza la situación desde una
lógica incremental: si existe capacidad ociosa y el pedido no canibaliza
ventas actuales, el costo relevante puede ser principalmente el costo
adicional de materiales y operación directa.

El pedido puede ser rentable aunque el precio sea menor, porque usa
capacidad que de otra manera estaría ociosa. Además, puede abrir nuevos
mercados.

\hypertarget{concepto-operacional-11}{%
\subsubsection{Concepto operacional}\label{concepto-operacional-11}}

Cuando hay capacidad ociosa en recursos no restringidos o capacidad
liberada del sistema, una venta incremental puede ser conveniente si
aumenta throughput más que el gasto operativo incremental y no desplaza
ventas más rentables.

Esto se conecta con la idea de que la restricción puede estar en el
mercado: si la planta puede producir más de lo que vende, el foco debe
moverse a generar demanda rentable.

\hypertarget{lo-evaluable-11}{%
\subsubsection{Lo evaluable}\label{lo-evaluable-11}}

La trampa típica es evaluar el pedido usando costo promedio o precio
normal, sin considerar capacidad ociosa ni contribución marginal.

\hypertarget{posible-vf-11}{%
\subsubsection{Posible V/F}\label{posible-vf-11}}

\textbf{Afirmación:} Un pedido con precio menor al precio local siempre
debe rechazarse.\\
\textbf{Respuesta:} Falso. Si usa capacidad ociosa, no desplaza ventas
actuales y cubre costos incrementales, puede aumentar throughput y ser
conveniente.

\hypertarget{capuxedtulo-39-nuevos-cuellos-de-botella-y-uso-luxf3gico-de-causa-efecto}{%
\subsection{Capítulo 39: Nuevos cuellos de botella y uso lógico de
causa-efecto}\label{capuxedtulo-39-nuevos-cuellos-de-botella-y-uso-luxf3gico-de-causa-efecto}}

\hypertarget{quuxe9-ocurre-12}{%
\subsubsection{Qué ocurre}\label{quuxe9-ocurre-12}}

Con los nuevos pedidos, parecen aparecer cuellos de botella por todos
lados. La situación muestra que la mejora no elimina los problemas para
siempre: cambia el sistema y aparecen nuevas restricciones.

Julie introduce a Alex a una forma de razonamiento lógico basada en
relaciones ``si\ldots{} entonces\ldots{}''. Esto ayuda a pensar efectos
secundarios antes de implementar una acción. El equipo analiza el
problema y concluye que debe proteger mejor los cuellos de botella con
inventario adecuado. Al aumentar el buffer, aumenta el tiempo de ciclo,
por lo que ventas debe ajustar promesas de entrega.

La decisión muestra una tensión real: bajar inventario es bueno, pero
bajarlo demasiado puede dejar sin trabajo al cuello de botella y
destruir throughput. El buffer no debe ser cero ni infinito; debe ser
suficiente para proteger el flujo.

\hypertarget{concepto-operacional-12}{%
\subsubsection{Concepto operacional}\label{concepto-operacional-12}}

El buffer es una protección contra variabilidad. Muy poco buffer puede
detener el cuello de botella. Demasiado buffer aumenta WIP y lead time.
La gestión correcta busca el tamaño razonable según variabilidad, ritmo
del cuello de botella y promesa comercial.

También aparece la idea de efectos secundarios: una mejora local puede
crear problemas si no se analizan sus consecuencias.

\hypertarget{lo-evaluable-12}{%
\subsubsection{Lo evaluable}\label{lo-evaluable-12}}

Puede salir una afirmación como ``siempre hay que minimizar inventario a
cero''. Eso es falso. TOC busca reducir inventario innecesario, pero
proteger el throughput.

\hypertarget{posible-vf-12}{%
\subsubsection{Posible V/F}\label{posible-vf-12}}

\textbf{Afirmación:} El buffer antes del cuello de botella debe
eliminarse para minimizar inventario.\\
\textbf{Respuesta:} Falso. El buffer protege al cuello de botella frente
a variabilidad. Eliminarlo puede detener el recurso que gobierna el
throughput.

\hypertarget{capuxedtulo-40-las-tres-preguntas-de-gestiuxf3n}{%
\subsection{Capítulo 40: Las tres preguntas de
gestión}\label{capuxedtulo-40-las-tres-preguntas-de-gestiuxf3n}}

\hypertarget{quuxe9-ocurre-13}{%
\subsubsection{Qué ocurre}\label{quuxe9-ocurre-13}}

Alex y Lou reflexionan sobre cómo administrar la división sin depender
de Jonah. Se dan cuenta de que los cinco pasos son útiles, pero
necesitan herramientas de pensamiento para aplicarlos a problemas
complejos. Surgen tres preguntas fundamentales:

\begin{enumerate}
\def\labelenumi{\arabic{enumi}.}
\tightlist
\item
  Qué cambiar.
\item
  Hacia qué cambiar.
\item
  Cómo provocar el cambio.
\end{enumerate}

Estas preguntas resumen el paso desde resolver problemas de planta a
liderar cambio organizacional. Alex entiende que debe volverse capaz de
diagnosticar, diseñar soluciones y conducir implementación sin crear
efectos secundarios graves.

\hypertarget{concepto-operacional-13}{%
\subsubsection{Concepto operacional}\label{concepto-operacional-13}}

La administración se vuelve un proceso de razonamiento. No se trata de
aplicar recetas, sino de construir una lógica de cambio:

\begin{itemize}
\tightlist
\item
  Identificar el problema raíz.
\item
  Diseñar una solución coherente con la meta.
\item
  Implementar considerando resistencia, secuencia y efectos secundarios.
\end{itemize}

\hypertarget{lo-evaluable-13}{%
\subsubsection{Lo evaluable}\label{lo-evaluable-13}}

Los capítulos finales conectan TOC con liderazgo y cambio
organizacional. La restricción puede ser física, de mercado, de política
o de pensamiento.

\hypertarget{posible-vf-13}{%
\subsubsection{Posible V/F}\label{posible-vf-13}}

\textbf{Afirmación:} El proceso de mejora continua termina cuando la
primera planta mejora.\\
\textbf{Respuesta:} Falso. La mejora continua no termina; al romper una
restricción aparece otra. El final del libro muestra la necesidad de
aplicar el método a toda la división.

\hypertarget{luxednea-narrativa-de-los-capuxedtulos-27-al-40}{%
\subsection{Línea narrativa de los capítulos 27 al
40}\label{luxednea-narrativa-de-los-capuxedtulos-27-al-40}}

Primero, la planta mejora pero Peach exige un 15\% adicional. Segundo,
Jonah propone reducir lotes para bajar inventario y lead time. Tercero,
ventas encuentra oportunidades que antes la planta no podía cumplir.
Cuarto, los reportes contables tradicionales empiezan a ser
cuestionados. Quinto, Alex es promovido y debe escalar el método. Sexto,
el equipo transforma la experiencia en los cinco pasos de mejora
continua. Séptimo, descubren que la inercia y las políticas antiguas
pueden convertirse en restricciones. Octavo, aparece el mercado como
posible restricción. Noveno, se refuerza la necesidad de buffers
adecuados y pensamiento lógico. Décimo, Alex concluye que administrar
exige responder qué cambiar, hacia qué cambiar y cómo provocar el
cambio.

\hypertarget{materia-de-operaciones-conectada-con-estos-capuxedtulos}{%
\subsection{Materia de operaciones conectada con estos
capítulos}\label{materia-de-operaciones-conectada-con-estos-capuxedtulos}}

\hypertarget{throughput-antes-que-eficiencia-local}{%
\subsubsection{Throughput antes que eficiencia
local}\label{throughput-antes-que-eficiencia-local}}

En estos capítulos se ve que una máquina ocupada no necesariamente
agrega valor. Si produce algo que no se vende o que no necesita el
cuello de botella, solo aumenta inventario. Esto es central para
Verdadero/Falso.

\hypertarget{lotes-muxe1s-pequeuxf1os}{%
\subsubsection{Lotes más pequeños}\label{lotes-muxe1s-pequeuxf1os}}

Reducir lotes puede:

\begin{itemize}
\tightlist
\item
  Reducir WIP.
\item
  Reducir lead time.
\item
  Mejorar flexibilidad.
\item
  Permitir entregas parciales.
\item
  Aumentar capacidad de respuesta comercial.
\end{itemize}

Pero hay que distinguir:

\begin{itemize}
\tightlist
\item
  En cuello de botella, demasiados setups pueden ser graves porque
  consumen capacidad crítica.
\item
  En no cuello de botella, más setups pueden ser aceptables si hay
  capacidad ociosa.
\end{itemize}

\hypertarget{restricciuxf3n-de-mercado}{%
\subsubsection{Restricción de mercado}\label{restricciuxf3n-de-mercado}}

Cuando la planta ya puede producir más, el problema pasa a ser vender
más. En ese caso, la restricción no está necesariamente en una máquina,
sino en la demanda.

\hypertarget{contabilidad-tradicional-versus-toc}{%
\subsubsection{Contabilidad tradicional versus
TOC}\label{contabilidad-tradicional-versus-toc}}

La contabilidad tradicional puede castigar:

\begin{itemize}
\tightlist
\item
  Menor eficiencia local.
\item
  Más setups.
\item
  Menor absorción de costos.
\end{itemize}

Pero TOC pregunta:

\begin{itemize}
\tightlist
\item
  ¿Sube el throughput?
\item
  ¿Baja el inventario?
\item
  ¿Se controla el gasto operativo?
\item
  ¿Mejora el flujo?
\end{itemize}

\hypertarget{inercia}{%
\subsubsection{Inercia}\label{inercia}}

La inercia aparece cuando una organización sigue usando una regla que
fue útil antes, pero que ya no corresponde a la restricción actual.

\hypertarget{buffer-correcto}{%
\subsubsection{Buffer correcto}\label{buffer-correcto}}

El buffer antes del cuello de botella no es desperdicio automáticamente.
Es protección del throughput. El error es creer que todo inventario es
malo o que todo inventario es bueno. El inventario correcto depende de
su función.

\hypertarget{banco-de-verdaderofalso-probable}{%
\subsection{Banco de Verdadero/Falso
probable}\label{banco-de-verdaderofalso-probable}}

\hypertarget{section}{%
\subsubsection{1}\label{section}}

\textbf{Afirmación:} En los capítulos finales, la mejora de la planta
elimina definitivamente las restricciones.\\
\textbf{Respuesta:} Falso. La mejora rompe una restricción, pero puede
aparecer otra. Por eso el quinto paso exige volver al paso 1.

\hypertarget{section-1}{%
\subsubsection{2}\label{section-1}}

\textbf{Afirmación:} Cuando la planta logra producir más de lo que
vende, la restricción puede moverse al mercado.\\
\textbf{Respuesta:} Verdadero. Si hay capacidad productiva disponible
pero no hay demanda suficiente, el throughput queda limitado por ventas.

\hypertarget{section-2}{%
\subsubsection{3}\label{section-2}}

\textbf{Afirmación:} Reducir lotes puede disminuir inventario en proceso
y tiempo de flujo.\\
\textbf{Respuesta:} Verdadero. Lotes más pequeños permiten que las
piezas avancen antes, reducen WIP y pueden acortar lead time.

\hypertarget{section-3}{%
\subsubsection{4}\label{section-3}}

\textbf{Afirmación:} Si reducir lotes baja la eficiencia local de una
máquina no cuello de botella, entonces necesariamente empeora el
sistema.\\
\textbf{Respuesta:} Falso. Un no cuello de botella puede tener capacidad
ociosa. Lo relevante es si aumenta throughput y reduce inventario, no si
la eficiencia local baja.

\hypertarget{section-4}{%
\subsubsection{5}\label{section-4}}

\textbf{Afirmación:} Un pedido con bajo precio siempre destruye valor.\\
\textbf{Respuesta:} Falso. Si usa capacidad ociosa, no desplaza pedidos
mejores y cubre costos incrementales, puede aumentar throughput.

\hypertarget{section-5}{%
\subsubsection{6}\label{section-5}}

\textbf{Afirmación:} La inercia puede transformarse en una
restricción.\\
\textbf{Respuesta:} Verdadero. Una política creada para una restricción
antigua puede volverse dañina si se mantiene cuando el sistema cambió.

\hypertarget{section-6}{%
\subsubsection{7}\label{section-6}}

\textbf{Afirmación:} Los cinco pasos de mejora continua son una receta
solo para manufactura física.\\
\textbf{Respuesta:} Falso. En los capítulos finales se aplican a la
división y a problemas de política, mercado y gestión.

\hypertarget{section-7}{%
\subsubsection{8}\label{section-7}}

\textbf{Afirmación:} El objetivo del buffer antes del cuello de botella
es proteger el throughput.\\
\textbf{Respuesta:} Verdadero. El buffer evita que el cuello de botella
quede sin trabajo por variabilidad aguas arriba.

\hypertarget{section-8}{%
\subsubsection{9}\label{section-8}}

\textbf{Afirmación:} El buffer óptimo siempre es cero porque todo
inventario es desperdicio.\\
\textbf{Respuesta:} Falso. Inventario excesivo es malo, pero un buffer
razonable antes del cuello de botella protege el flujo.

\hypertarget{section-9}{%
\subsubsection{10}\label{section-9}}

\textbf{Afirmación:} La contabilidad tradicional puede mostrar como
negativo un cambio que operacionalmente mejora throughput.\\
\textbf{Respuesta:} Verdadero. Indicadores basados en eficiencia local o
absorción de costos pueden contradecir la lógica de throughput.

\hypertarget{section-10}{%
\subsubsection{11}\label{section-10}}

\textbf{Afirmación:} Las tres preguntas finales son: qué cambiar, hacia
qué cambiar y cómo provocar el cambio.\\
\textbf{Respuesta:} Verdadero. Representan la transición hacia procesos
de pensamiento para administrar sistemas complejos.

\hypertarget{section-11}{%
\subsubsection{12}\label{section-11}}

\textbf{Afirmación:} Subordinar significa que el cuello de botella debe
adaptarse al ritmo de los no cuellos de botella.\\
\textbf{Respuesta:} Falso. Es al revés: los no cuellos de botella se
subordinan al ritmo y necesidad del cuello de botella.

\hypertarget{section-12}{%
\subsubsection{13}\label{section-12}}

\textbf{Afirmación:} El éxito con el pedido grande muestra que
operaciones y ventas deben coordinarse.\\
\textbf{Respuesta:} Verdadero. La planta solo convierte capacidad en
throughput si ventas consigue demanda compatible con el flujo real.

\hypertarget{section-13}{%
\subsubsection{14}\label{section-13}}

\textbf{Afirmación:} Las reglas que funcionaron en la planta de Alex
deben copiarse literalmente a toda la división.\\
\textbf{Respuesta:} Falso. Se debe copiar el método de razonamiento, no
necesariamente cada herramienta concreta.

\hypertarget{section-14}{%
\subsubsection{15}\label{section-14}}

\textbf{Afirmación:} Una hora ahorrada en un recurso no cuello de
botella siempre equivale a una hora ganada para el sistema.\\
\textbf{Respuesta:} Falso. Si el recurso no limita el sistema, ahorrar
tiempo allí puede no aumentar throughput.

\hypertarget{preguntas-abiertas-probables-y-respuesta-esperada}{%
\subsection{Preguntas abiertas probables y respuesta
esperada}\label{preguntas-abiertas-probables-y-respuesta-esperada}}

\hypertarget{explique-por-quuxe9-reducir-lotes-puede-mejorar-el-desempeuxf1o-aunque-aumente-setups}{%
\subsubsection{Explique por qué reducir lotes puede mejorar el desempeño
aunque aumente
setups}\label{explique-por-quuxe9-reducir-lotes-puede-mejorar-el-desempeuxf1o-aunque-aumente-setups}}

Reducir lotes puede bajar el inventario en proceso y el tiempo de flujo.
Si el aumento de setups ocurre en recursos no cuello de botella, no
necesariamente reduce el throughput porque esos recursos tienen
capacidad ociosa. La mejora relevante es que el producto avanza más
rápido, se cumplen pedidos antes y se puede ofrecer menor lead time al
mercado.

\hypertarget{explique-quuxe9-significa-que-la-restricciuxf3n-se-mueva-al-mercado}{%
\subsubsection{Explique qué significa que la restricción se mueva al
mercado}\label{explique-quuxe9-significa-que-la-restricciuxf3n-se-mueva-al-mercado}}

Significa que la planta ya no está limitada principalmente por su
capacidad interna, sino por la cantidad de pedidos rentables
disponibles. La capacidad productiva existe, pero si no hay demanda, no
se genera throughput. En ese escenario, ventas y estrategia comercial
pasan a ser parte central de la mejora operacional.

\hypertarget{explique-por-quuxe9-la-inercia-puede-ser-una-restricciuxf3n}{%
\subsubsection{Explique por qué la inercia puede ser una
restricción}\label{explique-por-quuxe9-la-inercia-puede-ser-una-restricciuxf3n}}

La inercia ocurre cuando se mantiene una política antigua por costumbre.
Una regla que fue correcta para una restricción anterior puede impedir
ver la nueva restricción. Por eso el quinto paso no solo dice volver al
paso 1, sino evitar que la inercia se transforme en la nueva limitación
del sistema.

\hypertarget{explique-las-tres-preguntas-finales-de-gestiuxf3n}{%
\subsubsection{Explique las tres preguntas finales de
gestión}\label{explique-las-tres-preguntas-finales-de-gestiuxf3n}}

Qué cambiar significa identificar el problema raíz o restricción. Hacia
qué cambiar significa diseñar una solución que acerque al sistema a su
meta. Cómo provocar el cambio significa implementar la solución
considerando secuencia, resistencia organizacional y efectos
secundarios.

\hypertarget{mapa-mental-final-para-escuchar-antes-de-la-prueba}{%
\subsection{Mapa mental final para escuchar antes de la
prueba}\label{mapa-mental-final-para-escuchar-antes-de-la-prueba}}

La planta mejora, pero Peach exige más. Alex entiende que necesita
ventas. Jonah propone reducir lotes para bajar inventario y lead time.
La planta logra prometer entregas rápidas. Aparece un pedido grande y el
equipo lo vuelve factible con entregas parciales. La contabilidad
tradicional no refleja bien la mejora, pero Lou muestra que el desempeño
real es fuerte. Alex es promovido. Ahora debe aplicar el método a toda
la división. El equipo formaliza los cinco pasos: identificar, explotar,
subordinar, elevar y volver a empezar. Luego descubren que la inercia
puede convertirse en restricción. Aparece capacidad disponible y Johnny
busca mercado. Un pedido europeo de menor precio puede convenir si usa
capacidad ociosa y aumenta throughput. Con más pedidos aparecen nuevas
tensiones y se ajustan buffers. Al final, Alex entiende que administrar
requiere pensar: qué cambiar, hacia qué cambiar y cómo provocar el
cambio.

\hypertarget{tabla-final-de-verdaderofalso-recopilados}{%
\subsection{Tabla final de Verdadero/Falso
recopilados}\label{tabla-final-de-verdaderofalso-recopilados}}

\begin{landscape}
\footnotesize
\setlength{\tabcolsep}{4pt}
\renewcommand{\arraystretch}{1.12}
\begin{longtable}{@{}
  >{\raggedleft\arraybackslash}p{0.035\linewidth}
  >{\RaggedRight\arraybackslash}p{0.445\linewidth}
  >{\centering\arraybackslash}p{0.045\linewidth}
  >{\RaggedRight\arraybackslash}p{0.395\linewidth}@{}}
\toprule
\textbf{Nº} & \textbf{Afirmación} & \textbf{V/F} & \textbf{Justificación corta} \\
\midrule
\endhead
\bottomrule
\endlastfoot
1 & En los capítulos finales, la mejora de la planta elimina definitivamente las restricciones. & F & Romper una restricción no termina el proceso; aparece otra y se vuelve al paso 1. \\
2 & Cuando la planta logra producir más de lo que vende, la restricción puede moverse al mercado. & V & Si falta demanda, el throughput queda limitado por ventas, no por capacidad interna. \\
3 & Reducir lotes puede disminuir inventario en proceso y tiempo de flujo. & V & Lotes pequeños bajan WIP y aceleran el avance de piezas por el sistema. \\
4 & Si reducir lotes baja la eficiencia local de una máquina no cuello de botella, entonces necesariamente empeora el sistema. & F & En no cuellos de botella puede haber capacidad ociosa; importa el efecto global, no la eficiencia local. \\
5 & Un pedido con bajo precio siempre destruye valor. & F & Puede convenir si usa capacidad ociosa, no desplaza mejores ventas y cubre costos incrementales. \\
6 & La inercia puede transformarse en una restricción. & V & Una regla antigua puede bloquear la mejora cuando la restricción cambia. \\
7 & Los cinco pasos de mejora continua son una receta solo para manufactura física. & F & En el final se aplican a división, mercado, políticas y gestión. \\
8 & El objetivo del buffer antes del cuello de botella es proteger el throughput. & V & Evita que el cuello de botella quede sin trabajo por variabilidad aguas arriba. \\
9 & El buffer óptimo siempre es cero porque todo inventario es desperdicio. & F & Un buffer razonable protege el flujo; el exceso sí aumenta WIP y lead time. \\
10 & La contabilidad tradicional puede mostrar como negativo un cambio que operacionalmente mejora throughput. & V & Métricas de eficiencia local o absorción pueden ocultar mejoras sistémicas. \\
11 & Las tres preguntas finales son: qué cambiar, hacia qué cambiar y cómo provocar el cambio. & V & Son la base de los procesos de pensamiento para administrar sistemas complejos. \\
12 & Subordinar significa que el cuello de botella debe adaptarse al ritmo de los no cuellos de botella. & F & Es al revés: los no cuellos de botella se subordinan al cuello de botella. \\
13 & El éxito con el pedido grande muestra que operaciones y ventas deben coordinarse. & V & La capacidad se convierte en throughput solo si ventas consigue demanda compatible con el flujo. \\
14 & Las reglas que funcionaron en la planta de Alex deben copiarse literalmente a toda la división. & F & Se replica el método de razonamiento, no necesariamente la herramienta exacta. \\
15 & Una hora ahorrada en un recurso no cuello de botella siempre equivale a una hora ganada para el sistema. & F & Si el recurso no limita el sistema, ahorrar tiempo allí puede ser un espejismo. \\
\end{longtable}
\end{landscape}

\hypertarget{resumen-general-final}{%
\subsection{Resumen general final}\label{resumen-general-final}}

Los capítulos 27 al 40 muestran la maduración completa de la Teoría de
Restricciones. Al inicio de esta parte, Alex ya logró mejoras
importantes en la planta, pero Bill Peach le exige más. Esa exigencia
obliga a mirar el sistema completo: no basta con producir mejor si no
hay demanda suficiente. Por eso aparece una idea clave para el examen:
la restricción puede moverse. Primero estuvo en la planta, luego puede
pasar al mercado, a las políticas internas, a las métricas contables o
incluso a la forma de pensar de la organización.

La reducción de lotes es una de las decisiones operacionales centrales.
A primera vista parece mala porque puede aumentar setups y bajar
eficiencia local. Sin embargo, desde TOC se evalúa globalmente: si se
reducen lotes en recursos no cuello de botella, se puede bajar WIP,
acortar lead time y ganar flexibilidad sin sacrificar throughput. Esto
permite aceptar pedidos grandes mediante entregas parciales y convertir
la rapidez operacional en ventaja comercial.

La contabilidad tradicional aparece como una fuente de conflicto. Puede
castigar decisiones correctas porque mira eficiencia local, utilización
o absorción de costos. TOC exige mirar throughput, inventario y gasto
operativo. Si aumenta el throughput real, baja inventario innecesario y
se mantiene controlado el gasto operativo, la empresa puede estar mejor
aunque algunos indicadores locales se vean peor.

El cierre del libro generaliza el aprendizaje. Alex deja de depender de
Jonah y entiende que administrar es pensar sistémicamente. Los cinco
pasos de mejora continua ordenan la acción: identificar, explotar,
subordinar, elevar y volver a empezar. El último paso es crucial porque
advierte contra la inercia: una política que antes ayudaba puede
transformarse en la nueva restricción si se mantiene después de que el
sistema cambió.

La idea final es que la gestión no consiste en copiar herramientas como
etiquetas, buffers o reglas de liberación. Consiste en saber
diagnosticar el sistema, encontrar la restricción actual, diseñar una
solución coherente con la meta y provocar el cambio sin generar efectos
secundarios mayores. Por eso las tres preguntas finales son tan
importantes: qué cambiar, hacia qué cambiar y cómo provocar el cambio.

\end{document}
